\documentclass[12pt]{article}
\usepackage[utf8]{inputenc}
\usepackage[T1]{fontenc}
\usepackage[margin=1in]{geometry}
\usepackage{setspace}
\usepackage{xcolor}
\usepackage[colorlinks=true, urlcolor=blue!60!black, linkcolor=blue!60!black]{hyperref}
\setlength{\parindent}{0em}
\setlength{\parskip}{0.6em}
\renewcommand{\baselinestretch}{1.0}
\begin{document}
\begin{center}
{\LARGE \textit{Milliken v.\ Bradley}, 418 U.S.\ 717 (1974)\par}
\vspace{0.5em}
{\large Case Brief --- Two-Page Edition\par}
\end{center}

\section{Memory Jogger}
No metropolitan school desegregation remedy across suburban lines without proof of an inter-district violation (\href{https://supreme.justia.com/cases/federal/us/418/717/}{Justia}).

\section{Detailed Case Facts}
Detroit's schools, about 63.6\% Black in a white metropolitan area, were found \textit{de jure} segregated by the Detroit Board and the State (including through Act 48, which nullified a desegregation plan). No proof was taken that any suburban district had segregated or that district lines were racially drawn.

\section{Procedural History}
Judge Roth rejected Detroit-only plans and ordered a remedy covering Detroit and 53 suburban districts; the en banc Sixth Circuit affirmed in principle.

\section{Judicial Votes}
Reversed, 5--4. Burger wrote for Stewart, Blackmun, Powell, and Rehnquist; Stewart concurred; Douglas, White, and Marshall dissented, each joined by the other dissenters.

\section{Holding}
The remedy is bounded by the violation. Cross-district relief requires that discriminatory acts were a substantial cause of inter-district segregation, or that lines were drawn by race.

\section{Analysis of Opinions}
\textbf{Majority:} The courts below sought metropolitan racial balance, which the Constitution does not require. District lines reflect deep traditions of local control, and consolidation would make the judge a super-superintendent; no cross-district effect was shown. \textbf{Stewart:} relief would be proper if lines were redrawn by race or state housing laws were used to discriminate. \textbf{Dissents:} Michigan, the entity the Fourteenth Amendment binds, cannot hide behind its subdivisions; a metropolitan plan was practicable; Marshall called the ruling ``a giant step backwards.''

\section{Examples: Future Applications}
\textbf{Same-Side:} gerrymandering within a city with neutral suburban borders confines relief to the city. \textbf{Opposite-Side:} a white town secedes to form its own district, allowing cross-district relief. \textbf{Fence-Sitter:} suburban racial zoning and redlining under Stewart's housing exception.

\section{Critique}
Defenders praise limits on judicial power; critics, including Myron and Gary Orfield, say the decision ignored state-sponsored housing segregation and, with \textit{Rodriguez}, entrenched metropolitan isolation.

\section{Key Quotations}
\begin{quote}
``No single tradition in public education is more deeply rooted than local control over the operation of schools; local autonomy has long been thought essential both to the maintenance of community concern and support for public schools and to quality of the educational process.''
\end{quote}
\noindent --- Chief Justice Burger, for the Court, 418 U.S.\ at 741--42.

\end{document}
